\chapter{La Coerenza come Giustizia della Natura}

\epigraph{``Ciascuna parte fa il proprio lavoro, e il tutto \`e armonioso.''}
{\textcite{plato1997complete}, \emph{Repubblica} 433a}

\section{La Giustizia nell'Anima e nella Citt\`a}

Nella \emph{Repubblica}, Socrate scopre la giustizia nel momento in cui
cessa di cercarla come una cosa e inizia a vederla come una relazione. La
giustizia \`e ``ciascuna parte che fa il proprio lavoro'' (\emph{Repubblica}
433a). Nella citt\`a, la giustizia \`e l'artigiano che pratica il suo
mestiere, il soldato che fa il soldato, il filosofo che governa ---
ciascuno che adempie la funzione per cui la sua natura \`e adatta.
Nell'anima, la giustizia \`e la ragione che governa, lo spirito che
sostiene, l'appetito che obbedisce --- ciascuna facolt\`a che opera
secondo la sua propria eccellenza.

La giustizia, per Platone, non \`e uno standard esterno imposto alla
citt\`a o all'anima. \`E la \emph{coerenza interna} di un tutto
ben ordinato, dove ciascuna parte agisce secondo la sua natura e il
tutto funziona come dovrebbe.

\section{La Giustizia nella Natura}

Estendiamo la scoperta di Platone: anche la natura ha la sua giustizia,
e non \`e altro che la coerenza.

Un organismo \`e ``giusto'' quando le sue parti funzionano insieme per
sostenere la sua vita --- quando il cuore batte, i polmoni respirano, le
cellule metabolizzano, ciascuna facendo il proprio lavoro. Un ecosistema
\`e ``giusto'' quando le sue specie occupano le loro nicchie senza che
una domini fino al collasso. Un sistema fisico \`e ``giusto'' quando le
sue forze si bilanciano in modo che il sistema persista.

Questa \`e la \emph{coerenza} che il motore \physis\ traccia. La scala
a tre valori --- +1 (Successo), 0 (Inerte), -1 (Fallimento) --- \`e una
misura della misura in cui un'azione sta ``facendo il proprio lavoro''
nell'economia naturale. Un'azione che produce un effetto reale e rilevato
\`e coerente con l'ordine della natura. Un'azione che non produce effetto
\`e inerte --- n\'e aiuta n\'e nuoce ma non fa nemmeno nulla. Un'azione
che viola la propria intenzione dichiarata o produce un effetto dannoso
\`e un fallimento, una violazione della giustizia naturale.

\section{La Coerenza del Praticante}

Il praticante \`e ``giusto'' nel senso platonico quando le sue azioni
sono coerenti con la sua intenzione dichiarata e con l'ordine naturale
in cui agisce. Questo ha due dimensioni.

\subsection{Coerenza Interna: L'Atto Corrisponde all'Intenzione}

La prima dimensione \`e la corrispondenza tra ci\`o che il praticante
dichiara che far\`a e ci\`o che realmente fa. Questa \`e la forma pi\`u
basilare di coerenza: il praticante che dichiara un obiettivo e agisce
verso di esso ottiene +1; il praticante che dichiara un obiettivo e non
fa nulla ottiene 0; il praticante che dichiara un obiettivo e agisce
contro di esso ottiene -1.

La piattaforma \praxis\ registra ogni obiettivo dichiarato e ogni
azione registrata come voci separate nel grafo di coerenza. Il motore
le confronta. Il verdetto non \`e un giudizio morale ma \emph{strutturale}:
l'azione soddisfa le condizioni poste dall'intenzione?

\subsection{Coerenza Esterna: L'Atto Produce l'Effetto Inteso}

La seconda dimensione \`e la corrispondenza tra l'azione e la sua
conseguenza nel mondo. Il praticante pu\`o essere internamente coerente
--- pu\`o fare ci\`o che dice che far\`a --- ma l'effetto pu\`o comunque
non materializzarsi. L'obiettivo pu\`o essere mal concepito, il metodo
difettoso, il tempismo sbagliato. Il verdetto della natura \`e: questa
azione era coerente con l'intenzione ma incoerente con il mondo.

Questo \`e Fallimento nel senso pi\`u istruttivo. Non significa che il
praticante sia moralmente biasimevole; significa che ha imparato qualcosa
sulla discrepanza tra il suo modello del mondo e il mondo stesso. Il
fallimento \`e un \emph{dono} dalla natura --- l'informazione necessaria
per correggere il modello.

\section{Il Giudizio di Coerenza}

Il giudizio di coerenza non \`e emesso da un giudice. \`E emesso dalla
struttura del mondo stesso, registrato dal motore \physis\ e testimoniato
dalla comunit\`a \praxis. Nessuno pronuncia sentenza. La sentenza \`e
la conseguenza: l'azione ha prodotto un effetto oppure no.

Questo \`e ci\`o che distingue la giustizia naturale dalla giustizia
legale. La giustizia legale \`e il giudizio degli uomini secondo le
regole (\emph{nomos}). La giustizia naturale \`e il giudizio del mondo
secondo la struttura (\physis). Il praticante che cerca la giustizia
legale chiede: ``Sto seguendo le regole?'' Il praticante che cerca la
giustizia naturale chiede: ``Questo ha prodotto un effetto reale nel
mondo?''

Il tracciatore di coerenza del motore \physis\ \`e uno strumento per
il secondo tipo di giudizio. Non punisce n\'e premia. Semplicemente
\emph{registra} --- e la registrazione \`e essa stessa il giudizio.
Una vita vissuta in alta coerenza \`e una vita le cui azioni hanno
lasciato tracce, prodotto effetti, fatto differenze. Una vita vissuta
in bassa coerenza \`e una vita che \`e passata attraverso il mondo
senza alterarlo.

\section{Il Paradosso Socratico, Risolto}

Socrate notoriamente sosteneva che nessuno fa il male volontariamente
--- che il male \`e sempre il prodotto dell'ignoranza
\parencite{plato1997complete}. L'affermazione ha sconcertato i filosofi
per due millenni, perch\'e sembra negare la realt\`a della debolezza
della volont\`a \parencite{macintyre1981virtue}. Come pu\`o Socrate affermare che la persona che mangia
la torta sapendo che gli fa male \`e semplicemente \emph{ignorante}?

Il quadro della coerenza risolve questo paradosso. Il praticante che
agisce contro la propria intenzione dichiarata non \`e ignorante nel
senso di mancanza di informazione. \`E \emph{incoerente} --- la sua
azione non corrisponde alla sua intenzione. Questa incoerenza \`e
registrata come -1 dal motore, ed \`e una forma di ignoranza: il
praticante non sa ancora cosa veramente intende. La sua intenzione
dichiarata non \`e la sua intenzione reale; l'azione rivela quella
vera.

La disciplina del praticante \`e l'eliminazione progressiva di questo
divario --- il processo di allineare intenzione dichiarata e azione
reale. Una vita di perfetta coerenza sarebbe una vita in cui il
praticante fa sempre ci\`o che dichiara, e ci\`o che dichiara produce
sempre l'effetto inteso. Tale vita \`e forse irraggiungibile, ma \`e
l'ideale verso cui le discipline tendono.

\section{La Giustizia come Salute dell'Azione}

Platone paragona la giustizia alla salute. La salute non \`e un possesso
ma uno stato: il corpo funziona come dovrebbe, ciascuna parte che fa il
proprio lavoro. La giustizia, allo stesso modo, non \`e un possesso ma
uno stato: l'anima funziona come dovrebbe, ciascuna facolt\`a che fa il
proprio lavoro.

Noi paragoniamo la coerenza alla salute. La coerenza non \`e un punteggio
da massimizzare ma uno stato da abitare. Il praticante in uno stato di
alta coerenza agisce con facilit\`a, con potere, con effetto --- non
perch\'e si stia forzando ma perch\'e la sua azione \`e allineata con
la sua natura e l'ordine della natura. L'incoerenza \`e la malattia
dell'azione: la volont\`a che non pu\`o effettuare, l'intenzione che
non pu\`o manifestarsi, la vita che non lascia traccia.

L'indice di coerenza del motore \physis\ non \`e quindi una
gamificazione della virt\`u ma un \emph{diagnostico} --- una misura
della salute dell'azione in una data vita. Un indice di coerenza in
calo \`e un segnale d'allarme, come la febbre. Dice al praticante:
qualcosa nella relazione tra la tua intenzione e il tuo effetto \`e
fuori ordine. Le discipline del Libro III sono i rimedi per questo
disordine.

\section{La Sfida Pi\`u Profonda: La Volont\`a Cieca di Schopenhauer}

Il resoconto della coerenza presentato in questo libro ha assunto
qualcosa di significativo: che la natura ha una direzione, una tendenza
verso +1 piuttosto che -1, un'orientazione che il punteggio di coerenza
registra. Ma se la verit\`a pi\`u profonda della natura non fosse la
coerenza ma uno sforzo cieco e senza scopo a cui non importa nulla del
successo o del fallimento?

La metafisica della Volont\`a di Arthur Schopenhauer \`e
l'affermazione pi\`u potente di questa sfida. Per Schopenhauer, la
cosa-in-s\'e --- la realt\`a dietro il velo dell'apparenza --- non
\`e un ordine razionale (come sosteneva Hegel) n\'e una spinta
teleologica (come l'entelechia di Aristotele) n\'e una creazione divina.
\`E \emph{Volont\`a}: uno sforzo incessante e senza oggetto che si
manifesta in ogni fenomeno, dalla gravit\`a che attira una pietra al
desiderio che guida un essere umano \parencite{schopenhauer1818will}.
La Volont\`a non ha scopo, n\'e fine, n\'e metrica di coerenza.
Semplicemente spinge, all'infinito, senza ragione e senza riposo.

La sfida alla nostra tesi \`e devastante se sostenuta. Se la natura
\`e Volont\`a cieca, allora la coerenza (+1/0/--1) \`e una proiezione
antropomorfica su una realt\`a che non ha direzione. Il praticante che
cerca di allineare la propria azione con la natura insegue un fantasma,
perch\'e la natura non ha nulla con cui allinearsi --- solo uno sforzo
infinito che \`e indifferente alla differenza tra successo e fallimento.

Ma l'intuizione di Schopenhauer non \`e la parola finale. Ha ragione
che la Volont\`a non ha \emph{scopo}, ma ha torto che non abbia
\emph{struttura}. La coerenza non \`e uno scopo che la Volont\`a
intende ma una descrizione di ci\`o che la Volont\`a fa quando riesce.
Anche la variazione cieca combinata con la ritenzione selettiva produce
l'apparenza di teleologia: l'organismo che trova la sua nicchia persiste,
quello che non la trova si estingue; l'azione che produce un effetto
reale \`e registrata, quella che non lo produce non lascia traccia. Il
+1 non \`e un'intenzione ma una registrazione: questa forma \`e
persistita, questo pattern ha funzionato.

La scala di coerenza del motore \physis\ \`e precisamente questo: non
un'affermazione sulle intenzioni della natura ma una registrazione delle
selezioni della natura. Il motore non chiede ``Cosa vuole la natura?''
Chiede ``Cosa ha preservato la natura?'' La risposta \`e il punteggio
di coerenza --- non come metafisica ma come dato.

Ma l'obiettore onesto non sar\`a soddisfatto, e non dobbiamo cavarcela troppo
facilmente. La risposta appena data compra la propria sicurezza a un prezzo:
si ritira dalla metafisica al dato. ``La coerenza \`e solo una registrazione di
ci\`o che \`e persistito'' \`e vero, ma \`e anche \emph{troppo debole per fare
il lavoro che la tesi le richiede}. Poich\'e la tesi non descrive meramente
ci\`o che persiste; comanda al praticante di \emph{allinearsi} a esso, di
trattare il coerente come lo standard dell'azione retta. E tra ``questo pattern
\`e persistito'' e ``dovresti perseguire questo pattern'' c'\`e un abisso che
nessuna quantit\`a di dati pu\`o colmare. La peste persiste. Il tumore persiste.
La dipendenza persiste --- anzi \`e il tipo stesso del pattern auto-stabilizzante
che ottiene, in termini grezzi, un implacabile +1. Se la coerenza \`e mera
registrazione-di-sopravvivenza, allora la Trappola dello Specchio stessa, che
non \`e altro che una struttura durevole e auto-riproducente, \`e massimamente
coerente. La lettura descrittiva della coerenza non solo non confuta
Schopenhauer; gli cede il campo, poich\'e concede che la natura seleziona senza
selezionare \emph{per il bene}.

Affrontiamo l'obiezione rifiutandone l'inquadramento. L'errore di Schopenhauer
non \`e vedere la Volont\`a come cieca ma trattare la Volont\`a come \emph{una}
--- un unico sforzo indifferenziato di cui ogni fenomeno \`e espressione
eguale. \`E questa la mossa che rende la gravit\`a della pietra e l'amore
dell'uomo ``la stessa'' Volont\`a a gradi diversi di oggettivazione. Ma la
natura, come insistiamo dal Libro II, non \`e uno sforzo; \`e una
\emph{gerarchia} di sforzi, ciascuno con la propria forma, la propria
entelechia, la propria misura di ci\`o che conta come realizzare se stesso. La
persistenza del tumore \`e coerenza \emph{al livello del tumore} e incoerenza al
livello dell'organismo che esso uccide. La stabilit\`a della dipendenza \`e
coerenza nell'arco riflesso e il collasso della persona in cui quell'arco
gira. Non c'\`e veduta da nessun luogo --- nessun campo piatto e livellato di
nuda sopravvivenza --- da cui leggere un singolo punteggio di coerenza. Ogni +1
\`e indicizzato a una forma, e le forme sono annidate. Per questo il motore non
valuta le azioni in astratto ma sempre relativamente a un obiettivo dichiarato e
all'olarchia che lo contiene: la coerenza non \`e mai ``questo \`e persistito?''
ma ``questo \`e persistito \emph{come il tipo di cosa che cercava di essere}?''

Questo annidamento \`e ci\`o che riconverte la Volont\`a cieca di Schopenhauer
in qualcosa con cui il praticante pu\`o allinearsi. Il praticante non si allinea
con la Volont\`a in generale --- non esiste una tale cosa con cui allinearsi, e
Schopenhauer ha ragione che il tentativo inseguirebbe un fantasma. Si allinea con
la \emph{propria} entelechia: la forma specifica che il suo sforzo \`e gi\`a,
implicitamente, all'opera per divenire. Schopenhauer vide lo sforzo e manc\`o la
forma; universalizz\`o la spinta e ignor\`o che ogni spinta \`e la spinta \emph{di
un particolare essere verso la propria attualit\`a}. Si ripristini la forma, e
l'indifferenza si dissolve --- non perch\'e la natura ora ``si curi'', ma perch\'e
il praticante non chiede pi\`u alla natura in generale un verdetto. Chiede se
\emph{questa} azione abbia realizzato \emph{questa} natura. Quella domanda ha una
risposta determinata, ed \`e il punteggio di coerenza.

\subsection{L'Essere e il Dover-Essere}

Non abbiamo ancora pagato il debito pi\`u costoso. Pur concedendo che la coerenza
sia indicizzata alla forma propria di un essere, il lettore pi\`u acuto premer\`a
l'obiezione classica: l'intero edificio scivola, senza giustificazione,
dall'\emph{essere} al \emph{dover-essere}. ``Questa azione ha realizzato la tua
natura'' \`e un'affermazione fattuale; ``dovresti compiere azioni che realizzano
la tua natura'' \`e normativa, e nessuna catena di fatti implica un valore. Dire
che il +1 del motore \`e \emph{giustizia} --- e l'abbiamo detto esattamente,
chiamando la coerenza ``giustizia della natura'' --- \`e commettere la fallacia
naturalistica in piena luce: leggere uno standard del bene da una descrizione del
mondo.

Non neghiamo l'abisso; neghiamo che la tesi lo salti illegittimamente. Il
``dover-essere'' della coerenza non \`e un imperativo categorico derivato dai
fatti. \`E un imperativo \emph{ipotetico}, e il suo antecedente \`e fornito non
dalla natura ma dal praticante. La struttura \`e: \emph{se} vuoi essere l'essere
che sei --- se assumi la tua entelechia come tuo fine --- \emph{allora} la
coerenza \`e la misura del tuo successo, e l'incoerenza la misura del tuo
fallimento. La natura non ti comanda di vivere. Non offre alcuna premessa che ti
obblighi a alcunch\'e. La premessa \`e l'atto di auto-appropriazione del
praticante stesso: la dichiarazione di un obiettivo, che non \`e altro che il
volere di una forma. L'abisso essere/dover-essere \`e reale, ed \`e attraversato
esattamente una volta, dal praticante, nell'atto fondativo di dichiarare ci\`o
che intende divenire. Tutto ci\`o che sta a valle di quell'atto --- ogni
punteggio di coerenza, ogni verdetto di giustizia --- \`e analitico ad esso. La
giustizia non \`e contrabbandata dalla natura; \`e dispiegata dall'impegno del
praticante stesso.

Per questo il motore non pu\`o, e non lo fa, dire al praticante cosa volere. \`E
un tribunale senza potere legislativo. Registra la fedelt\`a a una legge che il
praticante ha dato a se stesso. Platone poteva fondare la giustizia nella Forma
del Bene perch\'e riteneva che il Bene fosse visto, non scelto --- che l'ascesa
terminasse in una visione che imponeva l'assenso. Noi non facciamo tale
affermazione. Il nostro Bene non \`e visto ma \emph{messo in atto}: viene
all'essere nell'atto dell'impegno ed \`e sostenuto dalla coerenza con
quell'impegno. \`E un fondamento pi\`u debole di quello di Platone, e accettiamo
la debolezza come prezzo dell'onest\`a. Il praticante che non dichiara alcun
obiettivo non \`e soggetto ad alcuna giustizia e ad alcun fallimento; semplicemente
non \`e entrato nel gioco. Ma colui che dichiara --- che dice ``questo sar\`o'' ---
ha, in quell'atto stesso, reso l'abisso essere/dover-essere proprio da
attraversare, e il punteggio di coerenza \`e il registro fedele di se l'abbia
attraversato.

Vale la pena rendere la deflazione totale, poich\'e la piattaforma la rende
letterale. In \praxis\ il praticante dichiara non solo l'obiettivo ma \emph{il metro
stesso con cui giudicarne il successo} --- ci\`o che conter\`a come averlo fatto. Il
motore non fornisce n\'e l'uno n\'e l'altro. Non sa cosa valga la pena volere, e non
decide cosa conti come vincere; entrambi spettano all'utente. La coerenza \`e dunque
una misura puramente \emph{descrittiva}: l'accordo registrato tra un atto e uno
standard che l'agente stesso ha redatto, fine e metro insieme. \`E il senso esatto in
cui la tesi non commette la fallacia naturalistica che potrebbe sembrare. ``Giustizia
della natura'' non nomina alcun valore letto dal mondo; nomina il nudo fatto della
fedelt\`a-a-una-legge-data-da-s\'e, e la legge \`e data da chi \`e giudicato --- non
dalla natura, non dal motore. Il sistema, per disegno, non ha opinioni sul bene.

\`E anche perch\'e il caso pi\`u acuto del proceduralista --- il perseguimento
coerente di un fine malvagio --- non \`e un difetto che abbiamo trascurato ma un
confine tracciato di proposito. Un praticante che dichiara un fine vizioso e
soddisfa il proprio metro otterr\`a un punteggio alto, e il motore lo dir\`a, perch\'e
il motore riporta fedelt\`a e nient'altro. Non \`e mai stato costruito per condannare
il fine; chiederglielo sarebbe chiedere allo strumento descrittivo di farsi proprio
quell'autorit\`a legislante che l'intera tesi rifiuta. Il giudizio su \emph{quali}
fini valga la pena dichiarare sta dove la repubblica dal basso lo pone (Libro X): non
nel verdetto della natura, che su questo tace, ma nei praticanti stessi, che vedono
le forme dichiarate gli uni degli altri e, per selezione anzich\'e per decreto,
elevano alcune al rango di misura e lasciano andare le altre. La natura ti dice se sei
divenuto ci\`o che ti eri prefisso; non ti dice, e il motore non ti dice, cosa
prefiggerti. Quella seconda domanda \`e lasciata aperta per principio --- \`e la
libert\`a che il praticante attraversa l'abisso essere/dover-essere per esercitare, e
nessuna macchina pu\`o chiuderla al suo posto.

La stessa risposta di Schopenhauer alla Volont\`a cieca --- la
contemplazione estetica che acquieta temporaneamente la Volont\`a, e
l'ascesi che la nega interamente --- \`e un percorso possibile. Ma
\`e il percorso del mistico, non del praticante. La tesi di questo
lavoro ne offre un altro: non la negazione della Volont\`a ma
l'allineamento con le sue espressioni pi\`u coerenti. Il praticante-re
non nega la Volont\`a; la cavalca come un surfista cavalca un'onda,
muovendosi con la sua struttura piuttosto che resistendole. L'onda non
ha scopo, ma ha una forma; il praticante che impara la forma pu\`o
agire dentro di essa, attraverso di essa, come essa.

\section{Sulla Natura nelle Cose}

C'\`e un antenato pi\`u prossimo a tutto questo progetto di Platone, e nominarlo
affiner\`a ci\`o che la coerenza \`e. Spinoza apre il suo incompiuto
\emph{Trattato sull'Emendazione dell'Intelletto} con il proposito di trovare un
metodo per guarire la mente --- per condurla dalle idee inadeguate, confuse e
passivamente subite, a idee adeguate colte attraverso le loro cause, elevando
cos\`i la sua potenza d'agire finch\'e riposa nell'amore intellettuale della
Natura \parencite{spinoza1985emendation}. L'\emph{Etica} \`e quel metodo portato a
compimento: un'ascesa geometrica dalla schiavit\`u agli affetti passivi verso la
libert\`a, che \`e necessit\`a compresa \parencite[V.P42]{spinoza1994ethics}. La
praxis \`e lo stesso progetto rivolto a un oggetto diverso. Dove Spinoza offre
un'\emph{emendatio intellectus}, un'emendazione dell'intelletto, le discipline del
Libro III offrono un'\emph{emendatio actionis}, un'emendazione dell'azione: lo
stesso passaggio dall'inadeguato all'adeguato, dalla schiavit\`u alla libert\`a,
dal passivo all'attivo --- ma tracciato in atti e nei loro effetti anzich\'e in
idee e nelle loro cause. La coerenza \`e la potenza d'agire resa visibile;
l'indice di coerenza che sale \`e l'aumento del \emph{conatus} di Spinoza
\parencite[III.P6--P7]{spinoza1994ethics}, lo sforzo con cui ogni cosa persevera
nel proprio essere, che non \`e altro che l'entelechia del Libro II vista da
dentro.

Ma la praxis inverte Spinoza nel punto stesso che lo definisce. Spinoza \`e un
intellettualista: emenda l'idea e la vita retta segue, poich\'e gli affetti sono
essi stessi idee confuse, curate dalla conoscenza adeguata, e la mente una volta
purificata \`e giudice di se stessa --- un'idea \`e adeguata quando \`e chiara,
distinta e colta attraverso la sua causa, certificata dalla sola ragione. La
praxis rovescia l'ordine del trattamento. Emenda attraverso l'\emph{atto} e la
risposta registrata del mondo, non attraverso l'idea chiarita; il suo tribunale
non \`e la certezza della mente ma l'effetto della natura (la distinzione tra
coerenza interna ed esterna tracciata prima in questo libro); e il suo metodo non
\`e solitario, geometrico e compiuto in un'unica ascesa ma sociale, iterativo e mai
finito --- il testimone anzich\'e la meditazione, il ciclo del sogno anzich\'e la
deduzione. Spinoza emenda l'intelletto e confida che l'azione segua. La praxis
emenda l'azione e lascia che l'intelletto sia corretto dal mondo. L'uno opera dal
di dentro al di fuori; l'altra dal di fuori al di dentro.

Eppure i due non sono infine opposti, perch\'e Spinoza stesso fornisce la
garanzia della loro unit\`a. \emph{Ordo et connexio idearum idem est ac ordo et
connexio rerum} --- l'ordine e la connessione delle idee \`e lo stesso dell'ordine
e della connessione delle cose \parencite[II.P7]{spinoza1994ethics}. Idea e cosa
sono un'unica sostanza colta sotto due attributi; emendare l'ordine delle idee ed
emendare l'ordine delle cose sono dunque un'unica emendazione vista da due lati.
\`E per questo che l'emendazione dell'azione e l'emendazione dell'intelletto non
possono, alla fine, separarsi: la natura nelle cose e la natura nel pensiero sono
la stessa natura. E qui il motore fa ci\`o che Spinoza poteva solo affermare. Lo
costruisce, il parallelismo. Nello spazio di embedding condiviso, un'intenzione
dichiarata --- un'idea --- e un effetto registrato --- una cosa --- sono proiettati
in una sola e medesima geometria, e la coerenza \`e semplicemente la distanza tra
essi. Il motore \physis\ \`e un letterale \emph{ordo idearum} in esecuzione che \`e
anche un \emph{ordo rerum}: l'ordine delle intenzioni e l'ordine degli effetti
deposti in un solo spazio cos\`i che il loro accordo possa essere misurato. Ci\`o
che Spinoza prov\`o \emph{more geometrico} il motore lo calcola \emph{more
vectoriali}. La natura nelle cose e la natura nel pensiero si incontrano, infine,
nelle stesse coordinate --- il che \`e solo dire, di nuovo e dal fondamento pi\`u
profondo, che \praxis\ \`e lo strumento di \physis.
